\chapter*{Resumen}
\addcontentsline{toc}{chapter}{Resumen}

El aprendizaje profundo o Deep Learning ha transformado la inteligencia artificial, pero sus arquitecturas tradicionales están limitadas a datos con estructura euclídea regular (secuencias, rejillas, vóxeles). El \emph{Geometric Deep Learning} (GDL) emerge como un paradigma unificador que extiende estas técnicas a dominios no euclidianos (grafos, variedades, grupos) mediante la incorporación sistemática de simetrías geométricas. Esta tesis presenta un análisis matemático riguroso del GDL organizado en torno al marco de las «5~Gs»: Geometría, Grupos, Grafos, Geodésicas y Gauges.

El desarrollo parte de la formalización de las redes neuronales clásicas mediante teoría de la medida de Lebesgue, incluyendo una demostración completa del teorema de aproximación universal de Cybenko. El marco GDL se construye en torno a las cinco Gs: la \emph{Geometría} aporta la estructura de los dominios de datos mediante variedades riemannianas, espacios tangentes y curvatura; los \emph{Grupos} formalizan las simetrías a través de grupos y álgebras de Lie, teoría de representaciones (teoremas de Peter-Weyl y Schur) y el principio de equivarianza; los \emph{Grafos} discretizan dominios continuos mediante análisis espectral y el laplaciano de grafos; las \emph{Geodésicas} proporcionan métricas riemannianas, la aplicación exponencial y el transporte paralelo para la propagación de información; y los \emph{Gauges} introducen invariancias locales mediante fibrados principales y conexiones. Este aparato matemático se aplica a la formalización de redes convolucionales equivariantes (sobre grupos finitos, grupos de Lie compactos y no compactos, y variedades) y de mecanismos de atención geométrica, incluyendo Transformers equivariantes y atención hiperbólica.

A través de demostraciones formales y análisis detallados, esta tesis muestra cómo principios geométricos fundamentales (equivarianza, invarianza, localidad) unifican las redes convolucionales, las redes de grafos y los Transformers bajo un marco común de message passing geométrico. El resultado es una visión del GDL como marco matemático que no solo explica las arquitecturas existentes, sino que guía el diseño de nuevas.

\vspace{0.5cm}
\begin{flushleft}
	\textbf{Palabras clave}: Aprendizaje Profundo Geométrico, Teoría de Grupos, Equivarianza, Variedades Riemannianas, Teoría de Representaciones, Grupos de Lie, Redes Convolucionales, Redes de Grafos, Transformers, Message Passing
\end{flushleft}
